\chapter*{Forord}
\addcontentsline{toc}{chapter}{Forord}
\markboth{Forord}{Forord}
Velkommen til dette skrift.

Ved både 50 og 60 års jubilæet for \TKET blev der udarbejdet et jubilæumsskrift som tegnede et billede af \TKET. Det favnede både over hvordan foreningen har ændret sig gennem tiden, men også over hvordan ånden bag denne studenterforeningen eksisterede allerede fra de tidlige år.

Da det er lykkedes at sammensætte et jubilæumsskrift to gange i streg og \TKET er, om noget, traditionsbundet, skulle dette selvfølgelig gøres igen. I det der følger har vi taget udpluk af de gamle jubilæumsskrifter, selv nedskrevet historier fra det sidste årti, inddraget fortællinger fra de nuværende studerende, og gjort vores bedste for at male så levende et billede af nudagens \TK som muligt.

Vores håb er at gamle medlemmer kan læse dette skrift og nikke genkendende til det der bliver beskrevet, men også at nye studerende og folk der ikke på forhånd kender \TKET kan samle dette skrift op og lære \TKET lidt at kende.% samt få en fornemmelse af hvorfor så mange studerende har fundet et nyt hjem i \TKET baseret på et unikt sammenhold og en stærk ånd for frivilligt arbejde.

Hverdagen ved kammeret er en unik ting, med alle de skæve tiltag, spøjse traditioner og hyggelige stunder der opstår omkring denne studenterforening. Det er noget som de færreste af os havde forventet at vi ville få da vi startede, men de fleste af os er glad for at vi fik lov til at være en del af. \TKET har givet os bekendtskaber som vi tager med os resten af livet, og som ligger til grunde for det sammenhold og den entusiasme for fællesskabet som vi ser i vores medlemmer på tværs af årgange.

Det er takket været \TKETs medlemmer at vi har kunne sammensætte dette skrift, og vi har værdsat ethvert bidrag som vi har fået fra så vel gamle som nye medlemmer der har hjulpet med at male et billede af \TKET som det har set ud i løbet af det sidste årti.

\begin{flushright}
    J70 jubilæumsskriftsudvalget
\end{flushright}
